\documentclass[11pt]{article}
% Shared preamble for all solution files.
\usepackage[margin=1.1in]{geometry}
\usepackage{enumitem}
\usepackage{tikz}
\usepackage[hidelinks]{hyperref}
\usepackage{amsmath,amssymb,amsthm}
\newcommand{\fa}{\mathfrak a}
\newcommand{\fb}{\mathfrak b}
\newcommand{\fc}{\mathfrak c}
\newcommand{\fp}{\mathfrak p}
\newcommand{\fq}{\mathfrak q}
\newcommand{\fm}{\mathfrak m}
\newcommand{\fn}{\mathfrak n}
\newcommand{\fr}{\mathfrak r}
\newcommand{\Nil}{\mathfrak N}
\newcommand{\Jac}{\mathfrak R}
\newcommand{\rad}{r}
\newcommand{\Spec}{\operatorname{Spec}}
\newcommand{\Max}{\operatorname{Max}}
\newcommand{\Ann}{\operatorname{Ann}}
\newcommand{\Ker}{\operatorname{Ker}}
\newcommand{\Z}{\mathbf Z}
\newcommand{\Q}{\mathbf Q}
\newcommand{\R}{\mathbf R}
\newcommand{\C}{\mathbf C}
\newcommand{\ci}{\mathrm{c}}
\newcommand{\ex}{\mathrm{e}}

\setlength{\parindent}{0pt}
\setlength{\parskip}{0.55em}
% \exercise{label}{statement}  -- followed by the solution text
\newcommand{\exercise}[2]{\par\bigskip\noindent{\large\textbf{Exercise #1}}\par\smallskip
  \begin{quote}\small\textbf{Problem.}\ #2\end{quote}\par}
\newcommand{\pf}{\noindent\textbf{Solution.}\ }

\usepackage{needspace}
\AddToHook{cmd/section/before}{\Needspace{14\baselineskip}}
\let\chapterexercise\exercise
\renewcommand{\exercise}[2]{\Needspace{8\baselineskip}\chapterexercise{#1}{#2}}
\title{Liu, Chapter 7: Divisors and applications to curves\\[2pt]\large Solutions and corrections to the exercises}
\author{}
\date{}
\begin{document}
\maketitle
\textbf{Conventions.} The 52 labels follow the five exercise lists on printed pp.~263--266, 273--274, 283--284, 299--303 and 314--316 of the supplied text. Problems are restated with subparts in running prose. A regular element is a nonzerodivisor. Write $Q(A)$ for the total ring of fractions, $\mathcal K_X$ for the meromorphic-function sheaf, $L(D)=H^0(X,\mathcal O_X(D))$, and $l(D)=\dim_kL(D)$. For curves, $\chi_k(\mathcal F)=h^0(\mathcal F)-h^1(\mathcal F)$ and $p_a=1-\chi_k(\mathcal O_X)$. We use the chapter's Riemann--Roch theorem, Serre duality, the degree formula, the normalization sequence, and Riemann--Hurwitz, specifying their hypotheses when applying them. An ordinary multiple point means Liu's definition $\delta=m-1$, with $m$ the number of branches; it does not mean an arbitrary plane ordinary multiple point.

\textbf{Corrections and scope.} In 2.1 the UFD assertion is for a domain; a general normal Noetherian ring is a product of domains. Example~7.2.15 also has an omitted zero away from the cusp; see 2.4(c). The finite, nonflat version of 2.3(b) and the positivity criterion 2.5 are false; explicit nodal/cuspidal counterexamples are given. Exercise 2.3(a)'s unrestricted finite-map codimension assertion also fails; we give a finite birational pinching example, and prove the valid flat and normal-target versions. In 3.3(a) a nonzero section must be nonzero on every component. Exercise 3.4(a) has a twist-index shift. Exercise 3.5(d) needs quasi-compactness or a uniform pole bound. In 4.10 the Pell solution must be nonconstant. Exercise 4.11 uses the Weierstrass embedding with $o$ a flex, and 4.13(a) needs a hypersurface (for example Cohen--Macaulay) plane curve. In 5.7(c)--(e), integrality is valid but the proposed formulas for $\delta$ and the number of branches fail in general; we give counterexamples and valid criteria. These qualifications are stated again at the affected exercises.

\section*{7.1 Cartier divisors}

\exercise{1.1}{Let $X$ be locally Noetherian without embedded points, or reduced with finitely many irreducible components. For $U=\Spec A$ and nonnilpotent $f\in A$, let $a\in A$ be a zero divisor whose image in $A_f$ is regular. (a1) Show that $a$ is not nilpotent and that among the minimal primes $\fp_1,\ldots,\fp_n$, those containing $a$ are $\fp_1,\ldots,\fp_r$ with $0<r<n$. (a2) Choose $b\in\bigcap_{i>r}\fp_i\setminus\bigcup_{j\le r}\fp_j$; show some $m$ has $a^mb^m=0$ and $a^m+b^m$ regular. (a3) Show $a^{m-1}/(a^m+b^m)$ maps to $1/a$ and deduce $Q(A)_f\simeq Q(A_f)$. (b) Deduce quasi-coherence of $\mathcal K_X$ and $\mathcal K_X(U)=Q(A)$.}
\pf
Under either hypothesis the zero divisors are the union of finitely many minimal primes. Since $A_f\ne0$, a regular element in it is not nilpotent; this proves the first assertion. The zero divisor $a$ belongs to at least one minimal prime, and cannot belong to all, since then it would be nilpotent. Prime avoidance produces the stated $b$; the minimal primes are incomparable. The product $ab$ lies in every minimal prime, so is nilpotent. Choose $m$ with $(ab)^m=0$. At each minimal prime exactly one of $a,b$ lies in that prime, so their powers' sum lies in none and is regular.

In $A_f$, regularity of $a$ and $a^mb^m=0$ give $b^m=0$. Thus the displayed fraction maps to $1/a$. Every regular element of $A_f$ has the form $a/f^j$; if $a$ is already regular its inverse comes from $Q(A)$, and otherwise the construction above gives its inverse. This proves surjectivity of $Q(A)_f\to Q(A_f)$. For injectivity, an element $c/s$ vanishing after passage to $Q(A_f)$ has $c=0$ in $A_f$, since $A_f$ injects into its total fraction ring; some $f^Nc=0$, so it was already zero in $Q(A)_f$.

On the principal-open basis of $U$ the meromorphic presheaf is therefore exactly the localization presheaf of the $A$-module $Q(A)$. Its associated sheaf is $\widetilde{Q(A)}$. These affine identifications glue to $\mathcal K_X$, proving quasi-coherence and the asserted sections. Here the statement about the presheaf $\mathcal K'_X$ means this localization description on affine opens; a presheaf itself is not a quasi-coherent sheaf until sheafified. \hfill$\square$

\exercise{1.2}{If $X$ is locally Noetherian without embedded points, show that $X$ is reduced if and only if it is reduced at every generic point.}
\pf
The forward implication is immediate. Conversely the coherent nilradical is a subsheaf of $\mathcal O_X$, so all its associated points are associated points of $X$, hence generic. If it were nonzero, it would have an associated point on some affine open. At such a generic point its stalk is zero by hypothesis, a contradiction. Thus the nilradical vanishes. \hfill$\square$

\exercise{1.3}{If $X$ is locally Noetherian and has exactly one nonreduced local ring, at $x$, prove $x\in\operatorname{Ass}(\mathcal O_X)$.}
\pf
The nonzero coherent nilradical has support $\{x\}$. A nonzero finite module over a Noetherian ring has an associated prime, and any such prime lies in its support; hence its associated point is $x$. Since it is a submodule of $\mathcal O_X$, this is also an associated point of $X$. This argument is local in an affine neighborhood of $x$. \hfill$\square$

\exercise{1.4}{For Noetherian $X$ without embedded points, (a) prove $\mathcal K_X(V)=\prod_{\xi_i\in V}\mathcal K_{X,\xi_i}$, where the $\xi_i$ are generic points. (b) Prove $\mathcal K_X$ and $\mathcal K_X^*$ flasque and their positive-degree cohomology zero. (c) Deduce $\operatorname{CaCl}(X)\simeq\operatorname{Pic}(X)$.}
\pf
On an affine open, the primes of $Q(A)$ correspond to primes of $A$ consisting entirely of zero divisors. Prime avoidance and the absence of embedded primes show that these are precisely the minimal primes. Thus $Q(A)$ is a zero-dimensional Noetherian ring, hence Artinian, and its canonical product decomposition is
\[ Q(A)=\prod_{\fp\in\operatorname{Min}A}A_\fp. \]
These factors are exactly the generic meromorphic stalks. The decomposition commutes with restriction, which simply forgets absent generic points, and glues to prove (a). The product is finite, so it can also be written as a direct sum for the additive groups.

Restriction maps for both the rings and their unit groups are coordinate projections, hence surjective. Both sheaves are flasque and acyclic. The cohomology sequence of $1\to\mathcal O_X^*\to\mathcal K_X^*\to\mathcal K_X^*/\mathcal O_X^*\to1$ gives
\[ H^0(\mathcal K_X^*)\to\operatorname{Div}(X)\to\operatorname{Pic}(X)\to0. \]
Quotienting by principal divisors gives (c). \hfill$\square$

\exercise{1.5}{For a surjection $A\to B$ and a $B$-module $M$, prove $\operatorname{length}_A M=\operatorname{length}_B M$. Give a counterexample without surjectivity.}
\pf
The $A$-submodules and $B$-submodules of $M$ are exactly the same: every scalar in $B$ lifts to $A$. Thus the chains and their lengths coincide, including infinite length. For a field extension $k\subset L$ of finite degree $d>1$, $M=L$ has $L$-length $1$ but $k$-length $d$. \hfill$\square$

\exercise{1.6}{For a finite module $M$ over Noetherian local $(A,\fm)$, (a) characterize simple modules as $A/\fm$; (b) prove finite length equivalent to $\fm^rM=0$ for some $r$; (c) express its length as $\sum_i\dim_{A/\fm}\fm^iM/\fm^{i+1}M$; (d) if $A$ is a $k$-algebra, relate this to $\dim_k M$.}
\pf
(a) A nonzero element of a simple module generates it, giving $A/I$ with $I$ maximal, hence $I=\fm$. Conversely $A/\fm$ is simple. (b) A composition series of length $r$ has factors killed by $\fm$, so $\fm^rM=0$. Conversely such annihilation gives a finite filtration with finite-dimensional residue-field vector-space quotients. Refining each vector-space flag gives a composition series. Its length is the sum of these dimensions, proving (c).

Each simple factor has $k$-dimension $[A/\fm:k]$. Additivity along the composition series gives
\[ \dim_kM=\operatorname{length}_A(M)\,[A/\fm:k]. \]
For infinite residue degree the same statement is understood as a dimension/cardinality equality for nonzero $M$; the zero module is handled separately. \hfill$\square$

\exercise{1.7}{For finite $M$ define $e_A(a,M)=\operatorname{length}(M/aM)-\operatorname{length}(M[a])$ when both lengths are finite. (a) Show it is zero for finite-length $M$. (b) Show invariance after replacing $A$ by $A/I$ when $IM=0$. (c) For $0\to M'\to M\to M''\to0$, construct the two sequences with $N=(aM\cap M')/aM'$ and prove additivity, including finiteness if two pairs of lengths are finite.}
\pf
(a) In the endomorphism $a:M\to M$, the kernel and cokernel have equal length by the two image sequences. (b) The submodule lattices are unchanged by 1.5, so both defining lengths agree.

(c) The snake lemma for multiplication by $a$ gives
\[ 0\to M'[a]\to M[a]\to M''[a]\to M'/aM'\to M/aM\to M''/aM''\to0. \]
The image of the connecting map is $N$: a class in $M''[a]$ lifted to $m\in M$ maps to $am$ modulo $aM'$. Splitting at that image gives exactly
\[ 0\to N\to M'/aM'\to M/aM\to M''/aM''\to0,\qquad
0\to M'[a]\to M[a]\to M''[a]\to N\to0. \]
In any of the three choices of two finite pairs, these sequences express the remaining two modules as extensions of subquotients of finite-length modules. Thus all lengths are finite. Taking alternating sums gives $e_A(a,M)=e_A(a,M')+e_A(a,M'')$. \hfill$\square$

\exercise{1.8}{For a local homomorphism $A\to B$ put $e_{B/A}=\operatorname{length}_B(B/\fm_AB)$. (a) For an extension of DVRs prove $e_{B/A}=v_B(\pi_A)$. (b) If $M$ has finite $A$-length and $e_{B/A}<\infty$, prove $\operatorname{length}_B(M\otimes_AB)\le e_{B/A}\operatorname{length}_A M$, with equality if $B$ is flat.}
\pf
(a) In the target DVR, $\pi_A=u\pi_B^e$ for a unit $u$ and $e=v_B(\pi_A)$. The quotient by this element has length $e$. (b) Tensor a composition series of $M$ with $B$. Right exactness shows that each successive quotient in the induced filtration is a quotient of $B/\fm_AB$, so its length is at most $e_{B/A}$. Sum over the series. If $B$ is flat all the short exact sequences remain exact and the successive quotients are exactly $B/\fm_AB$, giving equality. \hfill$\square$

\exercise{1.9}{For quasi-projective $X$ over Noetherian affine $\Spec A$ and an ample invertible sheaf $\mathcal L$, prove that some $\mathcal L^m$ is $\mathcal O_X(D)$ for an effective Cartier divisor.}
\pf
Choose a sufficiently high power of $\mathcal L$ giving an immersion into projective space over $A$. The finite set $\operatorname{Ass}(X)$ maps to finitely many homogeneous primes not containing the irrelevant ideal. Homogeneous prime avoidance gives a positive-degree homogeneous form missing all these primes. One can see this without assuming an infinite base field: retain the inclusion-maximal primes, choose for each a homogeneous element lying in every other prime and not in that prime, raise them to a common positive degree, and add them. At any retained prime only its own summand survives. Homogeneous factors of positive degree outside the given prime allow the choices to have positive degree even when the original separating element has degree zero.

The corresponding section of a further power $\mathcal L^m$ vanishes at no associated point. In each local trivialization its coefficient is therefore a nonzerodivisor, by the associated-prime criterion. Its zero scheme is an effective Cartier divisor $D$ and its line bundle is $\mathcal L^m$. The zero divisor $D=0$ is allowed when the section is everywhere invertible. \hfill$\square$

\exercise{1.10}{Let $X/k$ be integral projective of dimension at least two and $D$ effective ample Cartier. (a) If $X$ is normal, prove $H^0(X,\mathcal O_X)\to H^0(nD,\mathcal O_{nD})$ is an isomorphism for large $n$, and $\operatorname{Supp}D$ connected. (b) Remove normality using normalization. (c) Give a reducible counterexample.}
\pf
(a) Use the negative-twist vanishing $H^1(X,\mathcal O_X(-nD))=0$ for $n\gg0$ on a normal projective variety of dimension at least two (the normal/S2 case of the vanishing theorem cited in the question, Hartshorne III.7.8). Also $H^0(X,\mathcal O_X(-nD))=0$: a nonzero such section would give a global regular function vanishing along $D$, whereas all global regular functions on integral projective $X$ form a field. The restriction sequence now gives the isomorphism. This field has no nontrivial idempotent, so $nD$, hence its underlying space $\operatorname{Supp}D$, is connected.

(b) Normalization $\pi:X'\to X$ is finite for varieties over a field. A finite pullback of an ample invertible sheaf is ample, and a finite scheme over a projective scheme is projective. Thus $\pi^*D$ satisfies (a). Its support maps surjectively onto $\operatorname{Supp}D$, proving connectedness.

(c) Take two planes in $\mathbf P^4_k$ meeting at exactly one point, and their reduced union $X$. A hyperplane avoiding that point cuts one line in each plane. Its Cartier hyperplane section is ample but is the disjoint union of those lines. The defining linear form avoids the generic points of this reduced scheme, so it is indeed a nonzerodivisor. \hfill$\square$

\exercise{1.11}{For locally Noetherian $X$, identify the direct limits of $\mathcal O_X(U)$ and $\mathcal O_X(U)^*$ over open $U\supseteq\operatorname{Ass}(X)$ with $H^0(X,\mathcal K_X)$ and its unit group, respectively.}
\pf
Proposition~7.1.15 gives $\mathcal K_X\simeq i_*\mathcal K_U$ for every such $U$. Thus a regular section there defines a unique global meromorphic section, and equality can be checked after restriction. Every meromorphic section has a local fraction $a/b$ with regular denominator. Its regularity locus is open and contains every associated point, since $b$ is invertible there. On that open it is an ordinary regular function, proving surjectivity of the first direct-limit map and its injectivity. If the meromorphic section is a unit, apply this to both it and its inverse and intersect their regularity loci. It is a regular unit there. This proves the group assertion as well. No uniform choice of a single denominator on $X$ is needed. \hfill$\square$

\exercise{1.12}{For coherent $\mathcal F$ on locally Noetherian $X$, define its associated points stalkwise. If $U\supseteq\operatorname{Ass}(\mathcal F)$, (a) prove $\mathcal F\to i_*(\mathcal F|_U)$ injective; (b) show a map $\phi:\mathcal F\to\mathcal G$ to a quasi-coherent sheaf is injective if its restriction to $U$ is.}
\pf
On an affine open a nonzero section of the kernel in (a) generates a nonzero finite submodule of the module of $\mathcal F$. This submodule has an associated prime, also associated to $\mathcal F$, in its support. Its support misses $U$, contradicting the hypothesis. This local argument proves (a). For (b), $\Ker\phi$ is a quasi-coherent submodule of coherent $\mathcal F$ and hence coherent on the Noetherian affine opens. Its associated points are among those of $\mathcal F$, yet its support misses $U$. The same argument makes the kernel zero. \hfill$\square$

\exercise{1.13}{For an invertible sheaf $\mathcal L$ on integral $X$ and a nonzero rational section $s$, (a) define $\operatorname{div}(s)$ from $s=e_if_i$ in local frames and identify its line bundle with $\mathcal L$; (b) recover principal divisors when $\mathcal L=\mathcal O_X$; (c) characterize effectiveness by regularity of $s$; (d) describe $\mathcal O_X(D)(U)$ as $\{f:\operatorname{div}(f)+D|_U\ge0\}\cup\{0\}$.}
\pf
In a frame $e_i$, the rational section is $e_if_i$ with $f_i\in K(X)^*$. On overlaps $f_i/f_j=e_j/e_i$ is a unit, so these functions define a Cartier divisor. The map $f_i^{-1}\mapsto e_i$ identifies its local generator with the frame of $\mathcal L$ and glues, proving (a). With the constant frame $1$ it is exactly the principal divisor of $s$, proving (b). Effectiveness says that each $f_i$ is regular, precisely that $s$ is regular, proving (c). Finally if $D$ has local equation $g_i$, membership of $f$ in $g_i^{-1}\mathcal O_X$ is equivalent to regularity of $fg_i$, which is exactly the effectiveness of $\operatorname{div}(f)+D$. Include zero separately, since a principal divisor requires a nonzero rational function. \hfill$\square$

\section*{7.2 Weil divisors}

\exercise{2.1}{For normal Noetherian $A$ with $\operatorname{Cl}(\Spec A)=0$, (a) show a height-one prime whose divisor is $(f)$ equals $fA$; (b) show irreducible elements generate prime ideals; (c) deduce unique factorization; (d) prove the converse vanishing of $\operatorname{Cl}$, $\operatorname{Pic}$ and $\operatorname{CaCl}$ for a Noetherian UFD.}
\pf
Assume first that $A$ is a domain. Normality gives $A=\bigcap_{\operatorname{ht}\fq=1}A_\fq$ in its fraction field. If $(f)=[\fp]$, all valuations of $f$ are nonnegative, so $f\in A$; its valuation is $1$ at $\fp$ and $0$ elsewhere. For $a\in\fp$, $a/f$ has nonnegative valuation everywhere, hence lies in $A$. Thus $\fp=fA$.

If $u$ is irreducible and $\fp$ is minimal over $(u)$, the principal ideal theorem gives height one. Part (a) gives $\fp=(f)$, and $u=fg$. Since $f$ is a nonunit, $g$ is a unit, so $(u)=\fp$. A Noetherian domain has factorization into irreducibles by the ascending chain condition on principal ideals, and primality of each irreducible gives uniqueness. Conversely every height-one prime in a Noetherian UFD is principal. Thus all Weil divisor classes vanish, and the injection of Cartier classes into Weil classes, together with the Cartier--Picard identification, gives the other vanishings.

If $A$ is not assumed a domain, its finitely many normal components are disjoint and $A$ is their product. Apply the argument to each factor. The conclusion is a product of UFDs, not a UFD in the usual domain-based definition; $k\times k$ shows why that qualification is necessary. \hfill$\square$

\exercise{2.2}{For number fields $L/K$, identify pushforward of a divisor on $\Spec\mathcal O_L$ with the norm of its fractional ideal and interpret the divisor push-pull formula.}
\pf
A divisor $\sum_{\mathfrak P}n_\mathfrak P[\mathfrak P]$ represents $\prod_\mathfrak P\mathfrak P^{n_\mathfrak P}$. Pushforward sends $[\mathfrak P]$ to $f(\mathfrak P/\mathfrak p)[\mathfrak p]$, exactly the formula $N_{L/K}(\mathfrak P)=\mathfrak p^{f(\mathfrak P/\mathfrak p)}$. Pullback represents extension of ideals, with exponents $e(\mathfrak P/\mathfrak p)$. Thus $f_*f^*D=[L:K]D$ says
\[ N_{L/K}(\mathfrak a\mathcal O_L)=\mathfrak a^{[L:K]},\qquad
\sum_{\mathfrak P\mid\mathfrak p}e(\mathfrak P/\mathfrak p)f(\mathfrak P/\mathfrak p)=[L:K]. \]
The analogous principal-divisor identity is the usual norm of elements. \hfill$\square$

\exercise{2.3}{Suppose $f:X\to Y$ is flat, or finite surjective between integral Noetherian schemes. For codimension-one $x$ put $y=f(x)$. (a) Bound $\dim\mathcal O_{Y,y}$, and prove finiteness of $e_{x/y}$ when that dimension is one. (b) The print asserts $\operatorname{mult}_x(f^*D)=e_{x/y}\operatorname{mult}_yD$ when $y$ has codimension one, and zero otherwise, and asks for the corresponding cycle formula.}
\pf
For the flat case put $A=\mathcal O_{Y,y}$, $B=\mathcal O_{X,x}$. Flat going-down gives $\dim A\le\dim B=1$. If $\dim A=1$, flat going-down again shows that a prime of $B$ containing $\fm_AB$ must be its maximal ideal: a nonmaximal prime in the one-dimensional local ring is minimal and cannot contract to a nonminimal prime. Thus $B/\fm_AB$ is Artinian and $e_{x/y}$ is finite.

If a local equation of an effective Cartier divisor at $y$ is a regular $a\in A$, its multiplicity is $\operatorname{length}_A(A/aA)$. Flatness and 1.8 give
\[ \operatorname{length}_B(B/aB)=e_{x/y}\operatorname{length}_A(A/aA). \]
Taking differences proves the same for arbitrary Cartier divisors. If $\dim A=0$, the Artinian local ring $A$ has only units as regular elements, so the pulled-back equation has multiplicity zero. Summing over codimension-one points proves exactly the requested cycle formula for flat $f$.

\textbf{The finite nonflat version of (b) is false.} Take the normalization $A=k[t^2,t^3]_{(t^2,t^3)}\subset B=k[t]_{(t)}$ and the divisor of $t^2$. Its multiplicity on $\Spec A$ is $2$, as $A/(t^2)$ has basis $1,t^3$; its pullback has multiplicity $2$. But $e_{B/A}=\operatorname{length}B/(t^2,t^3)=2$, so the asserted product is $4$.

For finite maps the quotient defining $e$ is Artinian whenever $x$ lies over $y$, since the fibre of a finite map is zero-dimensional. If $Y$ is normal, going-down for integral extensions and incomparability give equality of local heights, so $y$ has codimension one and $A$ is a DVR. The finite torsion-free algebra over this DVR is free, and its localizations are flat, so the displayed formula is valid in this corrected case as well. \textbf{The unrestricted finite-map assertion in (a) can itself fail.} Here is an explicit Noetherian example. Let $K=\Q(t_1,t_2,\ldots)$ and put
\[ R=K[u,v],\quad \fp=(u),\quad\fq=(u-1,v),\quad
B=(R\setminus(\fp\cup\fq))^{-1}R. \]
This Noetherian domain is semilocal with maximal ideals $\fp B,\fq B$ of heights one and two. Its residue fields are $K(v)$ and $K$. They are isomorphic as abstract fields via $\psi:K(v)\to K$, $v\mapsto t_1$, $t_i\mapsto t_{i+1}$. Let $I=\fp B\cap\fq B$ and define the subring
\[ A=\{b\in B:\psi(b\bmod\fp B)=b\bmod\fq B\}. \]
By the Chinese remainder theorem $A/I\simeq K$, diagonally inside $B/I\simeq K(v)\times K$. The element $e=1-u$ has residues $(1,0)$, and $B=A+Ae$: for any $b$, choose $a\in A$ with diagonal residue $\psi(b\bmod\fp B)-(b\bmod\fq B)$; then $b-ae\in A$. Thus $B$ is finite over $A$. The finite-subring theorem of Eakin--Nagata makes $A$ Noetherian. Integrality shows that $A$ is local with maximal ideal $I$ (both maximal ideals of $B$ contract to $I$), and $\dim A=\dim B=2$. The nonzero conductor $I$ also gives $Q(A)=Q(B)$, since $hB\subseteq A$ for any nonzero $h\in I$. Consequently $\Spec B\to\Spec A$ is finite birational and surjective, but its codimension-one point $\fp B$ maps to the dimension-two local point $I$. This violates precisely the claimed codimension equality without normality or an appropriate dimension hypothesis. \hfill$\square$

\exercise{2.4}{For a Cartier divisor $D$, (a) identify its support with $V(\mathcal O_X(-D))$ when effective; (b) compare its support with that of its Weil cycle, with equality for regular $X$ or effective $D$; (c) examine Example~7.2.15; (d) for a defined pullback of effective $E$, prove $\operatorname{Supp}[f^*E]=f^{-1}\operatorname{Supp}[E]$.}
\pf
(a) In a local equation $a$, the divisor is zero exactly where $a$ is a unit; this is the complement of $V(a)$. (b) At any codimension-one point where $D$ is zero its multiplicity is zero, giving $\operatorname{Supp}[D]\subseteq\operatorname{Supp}D$. If $D$ is effective, its support is the hypersurface cut out by a nonzerodivisor. Every component has codimension one by the principal ideal theorem (and cannot have codimension zero); its generic local length is positive. Thus the cycle has exactly that support. If $X$ is regular, the Cartier--Weil correspondence is locally an isomorphism, so vanishing of the cycle on an open is equivalent to vanishing of the Cartier divisor there.

(c) In the exact example of the text put $t=u^2,s=u^3$, so the normalization is $k[u]$ and the function $(t-s)/t$ is $1-u$. Let $p$ be the cusp $u=0$ and $q$ the smooth point $u=1$, namely $(s,t)=(1,1)$. At $p$ its order is zero, but it is not a unit of the cuspidal local ring, since $u\notin k[u^2,u^3]_{(u^2,u^3)}$. At $q$ it has a simple zero. Therefore
\[ \operatorname{Supp}D=\{p,q\},\qquad [D]=[q],\qquad\operatorname{Supp}[D]=\{q\}. \]
The printed claim that this function is a unit at every point other than $p$, and hence has zero cycle on the entire affine curve, omits $q$. On a neighborhood of $p$ excluding $q$ the advertised zero-cycle example is correct. A zero-cycle Cartier divisor on the local nodal ring is also described in 2.5 below.

(d) Pullback of an effective Cartier divisor, when defined, is cut out by the pulled-back regular equation. Its vanishing set is the inverse image of the original vanishing set. Apply the effective equality in (b) on both schemes. \hfill$\square$

\exercise{2.5}{For integral Noetherian $X$ with finite normalization $\pi:X'\to X$, the print asks to prove $[D]\ge0$ equivalent to $[\pi^*D]\ge0$, then equivalent to integrality of local equations of $D$ over $\mathcal O_X$.}
\pf
\textbf{Positivity of the cycle downstairs does not control each branch.} Over a field of characteristic different from two set
\[ A=k[t^2-1,t(t^2-1)]_{(t^2-1,t(t^2-1))}. \]
This is the integral local ring of a node. Its normalization is the semilocal ring
\[ B=k[t]_{\{(t-1),(t+1)\}}, \]
obtained by inverting polynomials nonzero at both $1$ and $-1$, with two residue fields $k$. Put $f=(t-1)^2/(t+1)$ and $D=\operatorname{div}(f)$ on $\Spec A$. The two normalization orders are $2,-1$. The length/order formula for a finite normalization of a curve says $\operatorname{ord}_A(f)=2+(-1)=1$, so $[D]\ge0$, whereas $[\pi^*D]$ has a negative coefficient. For clarity that order formula follows by choosing a common regular denominator, comparing the two multiplication sequences on $0\to A\to B\to B/A\to0$, and using 1.7(a) on the finite-length quotient.

The valid equivalence is
\[ [\pi^*D]\ge0\quad\Longleftrightarrow\quad
\pi^*\mathcal O_X(-D)\subseteq\mathcal O_{X'}\quad\Longleftrightarrow\quad
\mathcal O_X(-D)\subseteq\pi_*\mathcal O_{X'}. \]
Indeed a normal ring is the intersection of its height-one valuation rings, so nonnegative valuations mean that each local equation is regular upstairs. The last condition means exactly that it lies in the integral closure, hence is integral over the original affine ring. It implies $[D]\ge0$ by the finite normalization pushforward formula, but the converse fails as above. The function $(t-1)/(t+1)$ similarly has zero cycle downstairs but nonempty Cartier support. \hfill$\square$

\exercise{2.6}{For finite dominant $f:X\to Y$ with integral Noetherian schemes and normal $Y$, let $b\in K(X)^*$ and $a=N_{K(X)/K(Y)}b$. (a) Show $B=(f_*\mathcal O_X)_\xi$ finite free over $A=\mathcal O_{Y,\xi}$ at a height-one point. (b) Use elementary divisors over $A$ to compare norms and lengths. (c) Prove $f_*[\operatorname{div}(b)]=[\operatorname{div}(a)]$.}
\pf
(a) $A$ is a DVR and $B$ is finite torsion-free, hence free of rank $[K(X):K(Y)]$. (b) The elementary-divisor theorem gives bases for an inclusion $N\subset M$ of full-rank free modules with $N$ generated by $a_ie_i$ in a basis $e_i$ of $M$. For $b\in B\setminus\{0\}$ take $N=bB$, $M=B$. The determinant of multiplication by $b$ is a unit times $\prod_i a_i$: changing the domain basis can contribute a unit. Consequently
\[ v_A(N_{B/A}b)=\sum_i v_A(a_i)=\operatorname{length}_A(B/bB). \]
This also equals $\operatorname{length}_A(A/aA)$. The printed equality of the elements $a=\prod_i a_i$ should either include a unit or choose that unit in the invariant factors.

(c) Decompose the finite-length $B$-module $B/bB$ over its maximal ideals above $\fm_A$. Its $A$-length is
\[ \sum_{x\mapsto\xi}[k(x):k(\xi)]\operatorname{length}_{\mathcal O_{X,x}}(\mathcal O_{X,x}/b). \]
This is exactly the coefficient of $\xi$ in the pushed-forward principal divisor. The determinant valuation above is its coefficient in $\operatorname{div}(a)$. For a rational $b$, write it as a quotient of nonzero elements of $B$ and subtract the two identities; this also removes the temporary assumption $b\in B$. \hfill$\square$

\section*{7.3 Riemann--Roch theorem}

\exercise{3.1}{For projective $X/k$ and $D\in\operatorname{Div}(X)$, (a) show $D$ principal if and only if $L(D)$ contains a global meromorphic unit and $\deg D=0$. (b) If $X$ is integral, prove $D$ principal if and only if its pullback to any field extension $k'/k$ is principal.}
\pf
(a) If $D=\operatorname{div}(g)$, take $g^{-1}\in L(D)\cap H^0(\mathcal K_X^*)$; principal divisors on a projective curve have degree zero. Conversely such a unit $f$ gives an effective divisor $E=D+\operatorname{div}(f)$ of degree zero. A nonempty effective Cartier divisor on a curve has positive degree, the $k$-dimension of its finite coordinate algebra. Thus $E=0$ and $D$ is principal. This argument allows nonreduced curves.

(b) One direction follows by pullback. Conversely if $D_{k'}$ is principal, its degree is zero and $H^0(X_{k'},\mathcal O(D_{k'}))\ne0$. Flat cohomology base change gives $H^0(X,\mathcal O(D))\otimes_kk'\ne0$, so a nonzero section exists over $k$. Since $X$ is integral, this section is nonzero at its generic point and hence is a meromorphic unit in a rational trivialization of $\mathcal O(D)$. The invariant degree formula gives $\deg D=0$, and (a) applies. Geometric integrality of $X$ is not needed. \hfill$\square$

\exercise{3.2}{For normal projective integral $X$ over infinite $k$, distinct closed points $x_i$, and integers $n_i$, (a) show that $L(D+ny)\ne\bigcup_iL(D+ny-x_i)$ for large $n$, where $D=\sum n_ix_i$ and $y$ is another closed point. (b) Produce $f$ with $\operatorname{ord}_{x_i}f=n_i$.}
\pf
For large $n$, Riemann--Roch gives
\[ l(D+ny)-l(D+ny-x_i)=\deg x_i>0. \]
Indeed both divisors have degree greater than $2p_a-2$ after interpreting constants and duality over the base field. Thus every subspace on the right is proper. A finite-dimensional vector space over an infinite field is not a finite union of proper linear subspaces, proving (a). A closed point $y$ outside the finite prescribed set exists since an integral finite type curve has infinitely many closed points.

To get the signs in (b), apply (a) to $-D$. Choose $f\in L(-D+ny)$ in none of the subspaces $L(-D+ny-x_i)$. The first membership gives $\operatorname{ord}_{x_i}f\ge n_i$; exclusion of the $i$th subspace gives $\operatorname{ord}_{x_i}f<n_i+1$. These integral inequalities give equality. \hfill$\square$

\exercise{3.3}{For a rank-$r$ vector bundle $\mathcal E$ on a normal projective curve, (a) from a nonzero section $s$ form $\mathcal E'=(\mathcal E/s\mathcal O_X)/(\text{torsion})$ and show that its kernel is a line bundle; (b) construct such a line-subbundle quotient in general and deduce $\chi(\mathcal E)=\deg\det\mathcal E+r\chi(\mathcal O_X)$.}
\pf
First assume $X$ integral. A nonzero section is generically nonzero, so its image has generic rank one. Over a DVR every finite torsion-free module is free. Thus $\mathcal E'$ is locally free of rank $r-1$, and its kernel $\mathcal E''$ is locally free of rank one, as can be checked after splitting the surjection of free local modules. This proves (a).

For (b), twist by a sufficiently positive ample bundle to obtain a nonzero global section. Apply (a) and twist back. In
\[ 0\to\mathcal E''\to\mathcal E\to\mathcal E'\to0 \]
Euler characteristics add, and $\det\mathcal E=\mathcal E''\otimes\det\mathcal E'$. Induction on the rank, starting with line-bundle Riemann--Roch, gives the formula.

A normal curve may have disjoint components. In (a), a section which is zero on one component and nonzero on another does not have the claimed rank everywhere; one must assume it generically nonzero on every component. In (b) make the construction separately on each component and combine it; the Euler-characteristic formula still holds by additivity across components. \hfill$\square$

\exercise{3.4}{For projective $d$-dimensional $X/k$ with very ample $\mathcal O_X(1)$ and coherent $\mathcal F$, choose effective Cartier $D$ representing $\mathcal O_X(1)$, and put $Y=D$. (a) With $\mathcal G=\Ker(\mathcal F(-1)\to\mathcal F)$ show $\mathcal O_X(-D)\mathcal G=0$ and give the twist sequence. (b) Prove $\chi(\mathcal F(n))=P(n)$ for a polynomial of degree at most $d$, with $d!P\in\Z[T]$. (c) Identify it with $h^0$ for large $n$. (d) Compute its leading coefficient for an invertible sheaf on a curve.}
\pf
(a) Multiplication by a local equation of $D$ kills its own kernel, so $\mathcal G$ is a coherent sheaf on $Y$. Tensoring the original multiplication sequence by $\mathcal O_X(n+1)$ gives the correctly indexed sequence
\[ 0\to\mathcal G(n+1)\to\mathcal F(n)\to\mathcal F(n+1)\to(\mathcal F|_Y)(n+1)\to0. \]
The printed first term $\mathcal G(n)$ has an index shift. Therefore, writing $q(n)=\chi(\mathcal F(n))$,
\[ q(n+1)-q(n)=\chi((\mathcal F|_Y)(n+1))-\chi(\mathcal G(n+1)). \]
Both sheaves on the right live in dimension at most $d-1$.

(b) In dimension zero twists have the same Euler characteristic, the sum of the local lengths weighted by residue degrees. Inductively the displayed difference is a polynomial of degree at most $d-1$. Discrete antidifferentiation gives a polynomial $P$ of degree at most $d$ with that difference; choose its constant to equal $q(0)$, and recurrence in both directions gives $q(n)=P(n)$ for every integer $n$. An integer-valued rational polynomial of degree at most $d$ is an integral combination of $\binom{T}{j}$, $0\le j\le d$, by successive finite differences; hence $d!P\in\Z[T]$.

If the stipulated $D$ is unavailable for the chosen first twist over a finite field, make an infinite field extension first. Cohomology dimensions are unchanged by field extension, and over the infinite field one can choose a hyperplane avoiding the finite associated set. This proves the same polynomial statement over the original field.

(c) Serre vanishing makes all positive cohomology of $\mathcal F(n)$ zero for $n\gg0$. (d) Line-bundle Riemann--Roch on a projective curve gives
\[ P(T)=\deg\mathcal O_X(1)\,T+\deg\mathcal F+\chi(\mathcal O_X). \]
In particular the leading coefficient is independent of $\mathcal F$. \hfill$\square$

\exercise{3.5}{Let $X$ be locally Noetherian of dimension one, $U\supseteq\operatorname{Ass}(X)$ open, and $D$ effective Cartier with support $X\setminus U$. (a) Embed $\mathcal O_X(U)$ into $H^0(X,\mathcal K_X)$. (b) Prove $L(nD)\subseteq\mathcal O_X(U)$. (c) Locally at a boundary point clear a denominator by a large power of a local equation of $D$. (d) Prove $\mathcal O_X(U)=\bigcup_nL(nD)$.}
\pf
(a) Apply $\mathcal O_U\subset\mathcal K_U$ and the isomorphism $\mathcal K_X\simeq i_*\mathcal K_U$. (b) Restriction is injective by 1.12 applied to the invertible sheaf $\mathcal O_X(nD)$, whose associated points are those of $X$. On $U$ this sheaf is canonically $\mathcal O_U$.

(c) Write the meromorphic germ as $a/b$ with $b$ regular, and let $\pi$ cut out $D$. If $b$ is not a unit, all primes containing it are maximal in the one-dimensional local ring: a regular element avoids the minimal primes. Hence $\sqrt{bA}=\fm_A$, so $\pi^{n_x}\in bA$ for some $n_x$. Then $\pi^n(a/b)\in A$ for every $n\ge n_x$. If $b$ is a unit there is nothing to clear.

(d) If $X$ is Noetherian, the boundary is a finite set of closed points. Take the maximum of their bounds in (c); local membership everywhere gives $f\in L(nD)$. More generally quasi-compactness permits a finite subcover of the neighborhoods where such bounds work.

\textbf{The merely locally Noetherian assertion is false.} Take $X=\coprod_{i\ge1}\mathbf A^1_k$, $U=\coprod_{i\ge1}\mathbf G_{m,k}$ and let $D$ be the origin with multiplicity one in each component. The regular function whose $i$th component is $t^{-i}$ belongs to $\mathcal O_X(U)$, but to no $L(nD)$, since its pole orders are unbounded. \hfill$\square$

\section*{7.4 Algebraic curves}

\exercise{4.1}{For smooth geometrically connected projective genus-one $X$ with $o\in X(k)$, take bases $\{1,t\}$ of $L(2o)$ and $\{1,t,y\}$ of $L(3o)$. (a) Give the indicated bases through $L(5o)$. (b) Derive a Weierstrass relation in $L(6o)$ and normalize the coefficients of $y^2,t^3$ to one. (c) Identify the embedding by $|3o|$ with a smooth Weierstrass cubic.}
\pf
Riemann--Roch gives $l(no)=n$ for $n\ge1$, and $l(o)=1$. Thus $t$ has exact pole order $2$, $y$ exact pole order $3$, both only at $o$. The functions $1,t,y,t^2,ty$ have distinct pole orders $0,2,3,4,5$. They are independent and have the right cardinalities to give the requested bases.

There are seven displayed functions $1,t,y,t^2,ty,t^3,y^2$ in the six-dimensional space $L(6o)$, giving a relation
\[ b y^2+a_1ty+a_3y=a_0t^3+a_2t^2+a_4t+a_6. \]
If either $b$ or $a_0$ were zero, the other would also have to be zero by its unique order-six pole, contradicting the independence of the first five functions. Divide by $b$ and put $c=a_0/b$. The change $t'=ct$, $y'=cy$ makes the leading relation $y'^2=t'^3+\text{lower terms}$, without taking roots in $k$.

Finally $\deg(3o)=3>2g$, so its complete linear series is very ample. Its image in $\mathbf P^2$ has degree three, and the homogenized relation just found vanishes on it. The two cubic curves coincide, since the integral image has the same degree as this nonzero equation. Smoothness follows from the closed immersion. The unique point at infinity is $o=(0:1:0)$, a flex. \hfill$\square$

\exercise{4.2}{Show that a smooth geometrically connected projective curve of genus three is a plane quartic or becomes hyperelliptic over $\bar k$.}
\pf
Over $\bar k$, the canonical bundle has degree $4$, three sections, and no base point. If it is not very ample, there is a length-two divisor $p+q$ (allowing $p=q$) which it fails to separate. Riemann--Roch says
\[ l(K-p-q)=l(p+q), \]
so failure of separation gives $l(p+q)\ge2$. A genus-three curve cannot have a degree-one map to $\mathbf P^1$, so this degree-two linear system is basepoint free and defines a degree-two morphism. It is separable: a purely inseparable degree-two map would, by the Frobenius-twist description for smooth curves, force genus zero. Thus the curve is hyperelliptic.

Otherwise the canonical morphism is a closed immersion. Very ampleness descends under faithfully flat field extension, so the canonical series over $k$ embeds $X$ in $\mathbf P^2_k$. Its degree is $\deg K=4$, making the image a plane quartic. \hfill$\square$

\exercise{4.3}{Prove the converse of Lemma~7.4.3: a very ample effective divisor $D$ on a smooth connected projective curve over algebraically closed $k$ satisfies $l(D-p-q)<l(D-p)<l(D)$ for all $p,q$.}
\pf
A closed immersion by $|D|$ separates points and tangent vectors. Equivalently the restriction of its sections to any length-two subscheme $p+q$ is surjective; for distinct points this follows by a hyperplane separating their images, and for a double point by injectivity of the tangent map. Restriction to a single point is also surjective. The two kernels therefore have codimensions two and one in $L(D)$, respectively, giving both strict inequalities. \hfill$\square$

\exercise{4.4}{For a finite separable degree-$n$ map $f:X\to Y$ of normal projective curves, prove that $f$ is etale if and only if $\chi_k(\mathcal O_X)=n\chi_k(\mathcal O_Y)$.}
\pf
The different form of Riemann--Hurwitz (Thm.~7.4.16) gives
\[ 2p_a(X)-2=n(2p_a(Y)-2)+\deg_k R, \]
where $R$ is the effective different divisor. Its coefficient at $x$ is zero precisely when the extension of DVRs at $x$ is unramified, including separability of the residue extension. Finite maps of normal integral curves are flat, since finite torsion-free modules over DVRs are free. Thus the different vanishes precisely when $f$ is etale. Since every nonzero effective divisor has positive degree, the stated Euler-characteristic equality is equivalent to $R=0$. This argument does not require perfect $k$ or geometrically connected curves. \hfill$\square$

\exercise{4.5}{Let $f:X\to\mathbf P^1_k$ be finite separable with $X$ smooth geometrically connected. (a) If it is etale off infinity and tame there, prove it is an isomorphism; give a wild counterexample. (b) If it is etale off $0,\infty$ and tame there, show $K(X)=k(t)[h]$ with $h^n=\lambda t$, $n=\deg f$, $\lambda\in k^*$.}
\pf
Base change to $\bar k$, which preserves degree, genus and tame ramification. In (a), let $r$ be the number of geometric points above infinity. The tame Riemann--Hurwitz formula gives
\[ 2g(X)-2=-2n+(n-r)=-n-r. \]
Since $g\ge0$ and $n,r\ge1$, it forces $n=r=1$ and $g=0$. A finite degree-one map to the normal curve $\mathbf P^1$ is an isomorphism. In characteristic $p$, the map $\mathbf P^1_u\to\mathbf P^1_t$ with $t=u^p-u$ is a degree-$p$ counterexample when wild ramification at infinity is allowed: its derivative is $-1$ on the affine line.

(b) With $r_0,r_\infty$ geometric fibre cardinalities, the same calculation gives $2g-2=-r_0-r_\infty$. Hence $g=0$ and $r_0=r_\infty=1$. Tameness includes separable residue extensions, so the unique points above the two rational base points are $k$-rational; call them $P,Q$. Both have ramification index $n$, which is prime to the characteristic. Since $X$ has genus zero and a rational point, $X\simeq\mathbf P^1_k$. Choose a coordinate $h$ with divisor $P-Q$. Then $t$ has divisor $nP-nQ$, so $h^n/t=\lambda\in k^*$. Thus $K(X)=k(h)$ and $h^n=\lambda t$. If $k$ contains all $n$th roots of unity, $h\mapsto\zeta h$ gives the cyclic degree-$n$ covering. \hfill$\square$

\exercise{4.6}{For smooth geometrically connected projective $X/k$ of genus $g\ge2$, (a) show $\omega_X^3$ very ample; (b) construct a faithful representation of $\operatorname{Aut}_kX$ on its sections; (c) deduce finiteness over finite $k$.}
\pf
Its degree is $6g-6\ge2g+1$, so the usual degree criterion gives very ampleness. For $\sigma\in\operatorname{Aut}_kX$, use the natural differential pullback and define the action on sections by $(\sigma^{-1})^*$, so it is a homomorphism rather than an antihomomorphism. Its projectivization is the action on the embedding given by $\omega_X^3$. If it acts trivially on sections, it fixes that embedded curve pointwise and hence is the identity. Thus the representation is faithful. Over $\mathbf F_q$ its target is the finite group $\operatorname{GL}_{5g-5}(\mathbf F_q)$, since Riemann--Roch gives $h^0(\omega_X^3)=5g-5$. Therefore $\operatorname{Aut}_kX$ is finite. \hfill$\square$

\exercise{4.7}{For a separable hyperelliptic map $f:X\to\mathbf P^1$ of genus $g\ge2$, show that every rational point $x_0$ with $l(2x_0)\ge2$ is ramified. Assuming it unramified, (a) compute the poles of $h\pm\sigma(h)$ for $h\in L(2x_0)\setminus k$; (b) write $h=a(t)+b(t)y$ in the standard hyperelliptic equation with the fibre through $x_0$ placed at infinity; (c) derive $g\le1$.}
\pf
A nonconstant element of $L(2x_0)$ has exact pole order two at $x_0$: a single pole of order one would give a degree-one map to $\mathbf P^1$, contradicting the genus. If $x'_0=\sigma(x_0)\ne x_0$, $\sigma(h)$ has its only pole, also of order two, at $x'_0$. Thus at either point exactly one summand has a pole and no cancellation is possible. Both $h+\sigma(h)$ and $h-\sigma(h)$ have polar divisor $2x_0+2x'_0=2f^*(\infty)$ (in characteristic two these are the same function).

The affine chart in Prop.~7.4.24 is the entire inverse image of $\mathbf A^1_t$, with ring $k[t,y]$ and relation $y^2+Q(t)y=P(t)$. Since $h$ is regular there, $h=a(t)+b(t)y$ with $a,b\in k[t]$. Here $b\ne0$, since an invariant function has equal pole orders at the two unramified points, unlike $h$.

If the characteristic is not two, use $Q=0$. Unramified infinity forces $\deg P=2g+2$, and $y$ has pole order $g+1$ at both points. Hence $h-\sigma(h)=2b(t)y$ has pole order $\deg b+g+1$ there; (a) makes this equal to two, so $g\le1$. In characteristic two, unramified infinity forces $\deg Q=g+1$, and $h-\sigma(h)=b(t)Q(t)$ has pole order $\deg b+g+1$ at both points. The same conclusion follows. This contradicts $g\ge2$, proving the claim. \hfill$\square$

\exercise{4.8}{For smooth geometrically connected projective curves, (a) show a nonconstant map $X\to Y$ has $g(X)\ge g(Y)$; (b) if $g(X)\ge2$, show every finite separable self-map is an automorphism, and every finite self-map in characteristic $p$ is a relative Frobenius followed by an isomorphism from a Frobenius twist.}
\pf
Factor the map as a purely inseparable map followed by a separable map, using the maximal separable subextension of the function fields. The purely inseparable map from smooth $X$ identifies the intermediate normal curve with $X^{(p^r)}$ (Prop.~7.4.21); this twist is smooth and has the same genus. For a separable map of degree $n$, Riemann--Hurwitz gives $g(X)-1\ge n(g(Y)-1)$. If $g(Y)\ge1$ this implies the inequality in (a), and if $g(Y)=0$ it is automatic.

For a separable self-map with $g\ge2$, the same inequality forces $n\le1$, so it is a degree-one finite map between normal curves and is an isomorphism. In the general factorization the remaining separable map $X^{(p^r)}\to X$ is between curves of the same genus at least two, hence has degree one and is an isomorphism. All maps and twists here are over $k$; no identification of a twist with the original $k$-curve is presumed. \hfill$\square$

\exercise{4.9}{For smooth geometrically connected projective $X/k$ and nonempty open $U$, (a) identify divisors on $U$ with divisors on $X$ supported in $U$; (b) show the cokernel of $\alpha:\operatorname{Pic}^0(X)\to\operatorname{Pic}(U)$ is generated by any divisor of positive generator degree on $U$; (c) give an exact sequence with finitely generated kernel and cokernel; (d) over algebraically closed $k$, prove $\operatorname{Pic}(U)$ finitely generated exactly when $g=0$; (e) characterize smooth connected affine curves with finitely generated class group.}
\pf
(a) On a regular curve Cartier divisors are finite sums of closed points. A closed point of $U$ is closed in $X$, so extension by coefficient zero outside $U$ gives the identification.

(b) A line bundle on $U$ has a rational section and hence a divisor; extend its finite support to $X$. This also proves surjectivity of $\operatorname{Pic}(X)\to\operatorname{Pic}(U)$. Let $d_U>0$ generate the subgroup of degrees of divisors supported in $U$, and choose $D_0$ of degree $d_U$. For any such divisor $E$, $\deg E=md_U$ for some $m\in\Z$, so $E-mD_0$ has degree zero on $X$. Its restriction is in $\Imm\alpha$, proving the assertion.

(c) The kernel of restriction on $\operatorname{Pic}(X)$ is generated by the finitely many boundary points, since a trivialization on $U$ has divisor supported on the boundary. Its intersection $G$ with $\operatorname{Pic}^0(X)$ is therefore finitely generated. Part (b) makes $H=\operatorname{Coker}\alpha$ cyclic. Thus
\[ 0\to G\to\operatorname{Pic}^0(X)\xrightarrow{\alpha}\operatorname{Pic}(U)\to H\to0 \]
has the required properties, including $U=X$ with empty boundary.

(d) If $g\ge1$, choose a prime $\ell\ne\operatorname{char}k$. For every $m$, the group $\operatorname{Pic}^0(X)[\ell^m]$ is $(\Z/\ell^m\Z)^{2g}$, so $\operatorname{Pic}^0(X)$ has torsion of unbounded order and is not finitely generated. If $\operatorname{Pic}(U)$ were finitely generated, its subgroup $\Imm\alpha$ and the kernel $G$ would make $\operatorname{Pic}^0(X)$ finitely generated, a contradiction. For $g=0$, $X\simeq\mathbf P^1$ and $\operatorname{Pic}(U)$ is a quotient of $\Z$.

(e) Complete $\Spec A$ to its smooth projective model $X$. Regularity identifies $\operatorname{Cl}(A)$ with $\operatorname{Pic}(\Spec A)$. By (d) finite generation means $X\simeq\mathbf P^1$. Its affine open has a nonempty finite complement; choose one omitted point as infinity. It is then $\Spec k[x,1/P(x)]$ for a nonzero polynomial $P$ cutting out the remaining omitted points; a nonzero constant is allowed. Conversely each such ring is a localization of the PID $k[x]$, hence has trivial class group. \hfill$\square$

\exercise{4.10}{For separable even-degree $P(t)$ of degree at least two over algebraically closed $k$ of characteristic not two, let $X$ be the smooth completion of $y^2=P(t)$ and $x_1,x_2$ the points at infinity. Compare: (i) a polynomial Pell solution $A^2-PB^2=1$; (ii) $k[t,y]^*\ne k^*$; (iii) $x_1-x_2$ torsion in $\operatorname{Pic}^0(X)$.}
\pf
\textbf{In (i) require $B\ne0$, or equivalently a nonconstant solution.} Otherwise $A=1,B=0$ always works and carries no information. The affine ring $R=k[t,y]$ is normal because $P$ is separable and the characteristic is not two; its complement in $X$ is exactly $x_1,x_2$.

A solution gives $(A+By)(A-By)=1$, so $A+By$ is a unit, nonconstant if $B\ne0$. Conversely a nonconstant unit is uniquely $A+By$ with $A,B\in k[t]$. Its conjugate is a unit, and the norm $A^2-PB^2$ is a unit in $k[t]$, hence a nonzero constant $c$. Multiply both $A,B$ by $c^{-1/2}$, available over algebraically closed $k$, to get (i). A unit with $B=0$ would be constant, so the resulting solution is nontrivial.

Every such unit has divisor supported on the two points at infinity. Degree zero makes it $m(x_1-x_2)$; nonconstancy means $m\ne0$. Conversely a function with that principal divisor is regular and invertible on the affine complement, giving a nonconstant unit. This proves the corrected three-way equivalence. \hfill$\square$

\exercise{4.11}{For an elliptic plane cubic $E$ with origin $o=(0:1:0)$, suppose a plane curve $C$ meets $E$ only at $x\in E(k)$. Show $x-o$ is torsion.}
\pf
Use the standard Weierstrass embedding, so its line at infinity cuts $3o$ and $\mathcal O_E(1)=\mathcal O_E(3o)$. If $C$ has degree $d$ and no component $E$, its restriction is a nonzero section of $\mathcal O_E(d)$. Bezout, or the degree of that line bundle, says that its zero divisor has degree $3d$. Since it is supported at the rational point $x$, it is $3d\,x$. Thus $3d\,x\sim3d\,o$ and $3d(x-o)=0$ in $\operatorname{Pic}^0(E)$.

The flex condition on $o$ is necessary if ``elliptic plane cubic'' is interpreted as an arbitrary cubic with an arbitrary chosen origin. Over $\C$, take a flex $x$, let $C$ be its tangent line, and choose $o$ with $x-o$ nontorsion. A projective coordinate change puts $o$ at $(0:1:0)$, while $C\cap E=\{x\}$ still holds. Thus the coordinate label alone cannot replace the Weierstrass/flex hypothesis. \hfill$\square$

\exercise{4.12}{For smooth geometrically connected projective $X/\mathbf F_q$ with Jacobian $J$, (a) show $J(\mathbf F_q)$ finite; (b) for distinct rational points $x_1,\ldots,x_n$, $n\ge1$, let $F$ consist of nonzero functions whose divisors are supported on them. Show its order-vector image has finite index in $\{(r_i)\in\Z^n:\sum r_i=0\}$ and deduce $F\simeq k^*\oplus\Z^{n-1}$.}
\pf
(a) Every finite type scheme over a finite field has finitely many rational points, by a finite affine cover and its embeddings in affine spaces. Apply this to $J$.

(b) Let $\Lambda=\{(r_i):\sum r_i=0\}$. Sending a vector to $\mathcal O_X(\sum r_ix_i)$ gives $\Lambda\to\operatorname{Pic}^0(X)=J(k)$; the equality uses the rational point $x_1$. Its kernel is exactly the principal divisors with that support, hence $\Imm\varphi$. Part (a) makes this kernel finite index. It is a free abelian group of rank $n-1$. The kernel of $\varphi:F\to\Lambda$ is $k^*$, since a divisor-zero function is a global unit on $X$, and $H^0(X,\mathcal O_X)=k$. The resulting extension splits as an abstract abelian-group extension by choosing lifts of a basis of $\Imm\varphi$. This gives the stated noncanonical decomposition. \hfill$\square$

\exercise{4.13}{Call an affine curve plane if it is a closed subscheme of an open in $\mathbf A^2_k$. (a) Show its canonical sheaf trivial. (b) For smooth geometrically connected projective $X$ of genus $g\ge2$, bound by $(2g-2)^{2g}$ the rational points $x$ for which $X\setminus\{x\}$ is plane. (c) For hyperelliptic $X$, show these are precisely the rational Weierstrass points.}
\pf
(a) For a \emph{hypersurface} curve $C$ in an open $V\subset\mathbf A^2$, the conormal bundle and $\omega_V$ are both trivial. Indeed $\operatorname{Pic}(V)=0$, since $V$ is regular and its Weil class group is a quotient of the class group of the UFD $k[u,v]$. Adjunction then gives $\omega_C=\mathcal O_C$. The same applies to a pure Cohen--Macaulay plane curve, since it is locally a hypersurface in the regular surface.

\textbf{Embedded points must be excluded.} For $R=k[u,v]/(u^2,uv)$, the curve has an embedded point and is not l.c.i. Its canonical module computed as $\operatorname{Ext}^1_{k[u,v]}(R,k[u,v])$ is $k[u,v]/(u)$, not $R$: the resolution with relation $(-v,u)$ gives kernel generated by $(u,v)$ modulo $(u^2,uv)$. Thus the unrestricted wording of (a) does not imply its conclusion. All complements in (b),(c) are smooth, so this issue does not affect those parts.

(b) Write $n=2g-2$. Triviality of $\omega_X$ on $X\setminus\{x\}$ implies $K_X\sim n x$: a trivializing rational differential has divisor supported only at $x$, with coefficient its degree $n$. If there is one such $x_0$, every other gives $n(x-x_0)=0$ in $J(k)$. The Abel map $x\mapsto[x-x_0]$ is injective on rational points when $g\ge1$, since a principal divisor $x-y$ with $x\ne y$ would give a degree-one function. The number of geometric $n$-torsion points of a $g$-dimensional Jacobian is at most $n^{2g}$, also when the characteristic divides $n$ (use the prime-to-characteristic and characteristic-primary torsion formulas in Cor.~7.4.41). This gives the bound.

(c) If $x$ is Weierstrass, 4.7 makes it a ramified point of the hyperelliptic map. Place its image at infinity. The affine model of Prop.~7.4.24 is then exactly $X\setminus\{x\}$, a plane curve, proving one implication.

Conversely suppose the complement is plane, so $K_X\sim(2g-2)x$. For the hyperelliptic involution $\sigma$, the canonical class is $(g-1)$ times the hyperelliptic fibre class. If $x$ were unramified this would imply $(g-1)(x-\sigma x)$ principal. A corresponding function $h$ has degree $g-1$ and is not invariant under $\sigma$, since its pole is confined to one of these two points. Thus $(f,h)$ is birational onto its image in $\mathbf P^1\times\mathbf P^1$: the quadratic extension $K(X)/k(f)$ has no proper intermediate field. After algebraic closure this image has bidegree $(g-1,2)$, whose arithmetic genus is $g-2$, by the divisor restriction sequence on the product. The genus of its normalization $X$ is at most this, a contradiction. Therefore $x$ is ramified and $l(2x)\ge2$. \hfill$\square$

\exercise{4.14}{For a group scheme $G/S$, prove separatedness equivalent to the unit section being a closed immersion.}
\pf
If $G/S$ is separated, every section is a closed immersion, by base change of its diagonal. Conversely consider $d:G\times_SG\to G$, $(g,h)\mapsto gh^{-1}$. Scheme-theoretically $d^{-1}(\epsilon(S))$ is the diagonal: its functor of points consists exactly of pairs $g=h$. Thus the diagonal is a base change of the unit immersion, hence is closed when the latter is closed. This is separatedness. \hfill$\square$

\exercise{4.15}{For an algebraic group $G/k$ with unit $o$ and $x\in G(k)$, (a) construct translation carrying $o$ to $x$ and identify their local rings; (b) if $G$ is geometrically reduced, prove it smooth.}
\pf
Left multiplication $t_x:g\mapsto xg$ is a $k$-automorphism with inverse $t_{x^{-1}}$ and sends $o$ to $x$, giving the local-ring isomorphism. For (b) pass to $\bar k$. A reduced finite type scheme over a perfect field has a nonempty smooth open on each irreducible component. In particular there is a smooth closed point $a$. Every geometric closed point is a translate of $a$, so all closed points are smooth. The nonsmooth locus is closed, and a nonempty closed subset of a finite type scheme over a field has a closed point; hence that locus is empty. Smoothness descends under the faithfully flat extension $\bar k/k$. \hfill$\square$

\exercise{4.16}{Let geometrically reduced $G/k$ have morphisms $m,\operatorname{inv},\epsilon$ inducing an abstract group structure on $G(\bar k)$. Show they define an algebraic group.}
\pf
After passage to $\bar k$, every finite product of $G$ is reduced. The associativity maps on $G^3$, the two unit maps on $G$, and the two inverse composites on $G$ agree on every closed point by the group axioms on $G(\bar k)$. Apply Exercise~3.2.9: maps from a geometrically reduced variety which agree on all geometric points are equal. To recall its local argument, inverse images of any affine target open are the same, since their symmetric difference is constructible and a nonempty constructible set in a variety has a closed point. On that common open restrict further to affine source opens; the two maps to the affine target agree on functions because their difference vanishes at every closed point of a reduced affine algebra. This comparison does not require a separated target. Thus all group identities hold as morphisms after base change and descend by faithful flatness. It is essential to use geometric points; an abstract group law on $G(k)$ alone would not suffice over a finite field. \hfill$\square$

\exercise{4.17}{(a) Determine all $k$-automorphisms of $\mathbf A^1$. (b) For a group law with unit $0$, write its polynomial $P(Y,Z)$ and its values on rational points. (c) Prove the group is $\mathbf G_a$.}
\pf
(a) An automorphism sends $X$ to a polynomial $F$ with polynomial inverse $H$. Since $\deg(H\circ F)=\deg H\deg F=1$, both degrees are one: $F=aX+b$, $a\in k^*$.

(b) A morphism $\mathbf A^2\to\mathbf A^1$ is exactly a polynomial $P(Y,Z)\in k[Y,Z]$, and evaluation on a point gives $m(y,z)=P(y,z)$.

(c) Work over the infinite field $\bar k$. Each translation $Y\mapsto P(Y,z)$ is an automorphism, so has degree one in $Y$; since this holds for infinitely many $z$, all higher coefficients vanish. Similarly the degree in $Z$ is at most one. The identities $P(Y,0)=Y$ and $P(0,Z)=Z$ then give $P=Y+Z+cYZ$. If $c\ne0$, translation by $z=-1/c$ is constant, contradicting invertibility. Thus $c=0$. Equality of these polynomials descends to $k$, giving the standard additive law. \hfill$\square$

\exercise{4.18}{For a group law on $T=\Spec k[X,X^{-1}]$, normalize its unit to $1$. (a) Show the multiplication function is $aY^qZ^r$. (b) Prove $a=q=r=1$ in the appropriate senses and identify $T$ with $\mathbf G_m$.}
\pf
The units of the Laurent polynomial ring $k[Y^{\pm1},Z^{\pm1}]$ are exactly scalar monomials: considering extreme exponents in a product with its inverse forces its support to be a single exponent. Since the image of $X$ must be a unit, it is $aY^qZ^r$ with $a\in k^*$ and $q,r\in\Z$. The scheme-theoretic unit identities give $aY^q=Y$ and $aZ^r=Z$ in Laurent polynomial rings, forcing $a=1$ and $q=r=1$. Thus the group law is $YZ$. This argument avoids drawing conclusions about exponents merely from values on $k^*$, which would be insufficient over a finite field. \hfill$\square$

\section*{7.5 Singular curves and the structure of Pic0}

\exercise{5.1}{For reduced connected projective $X/k$, prove $k'=H^0(X,\mathcal O_X)$ is a finite field extension, give $X$ its natural $k'$-structure, and characterize $p_a(X)<0$ by $k'\ne k$ and $H^1(X,\mathcal O_X)=0$.}
\pf
The finite-dimensional $k$-algebra $k'$ is reduced since $X$ is reduced. It is therefore a product of finite fields. Connectedness excludes a nontrivial idempotent, so there is only one factor. The universal morphism $X\to\Spec H^0(X,\mathcal O_X)$ supplies its $k'$-structure; finite type over $k'$ follows from finite type over $k$.

Put $d=[k':k]$ and $h=\dim_{k'}H^1(X,\mathcal O_X)$. The cohomology group is naturally a $k'$-vector space, and
\[ p_a(X)=1-d+dh=1+d(h-1). \]
For integers $d\ge1,h\ge0$ this is negative exactly when $h=0$ and $d>1$. These are precisely the two stated conditions. \hfill$\square$

\exercise{5.2}{For projective $X/k$ with $H^0(X,\mathcal O_X)=k$ and any closed curve $Z\subset X$, prove $p_a(Z)\le p_a(X)$.}
\pf
The ideal sequence $0\to\mathcal I\to\mathcal O_X\to\mathcal O_Z\to0$ and vanishing of $H^2$ for coherent sheaves on a curve show $H^1(X,\mathcal O_X)\twoheadrightarrow H^1(Z,\mathcal O_Z)$. Also $H^0(Z,\mathcal O_Z)$ contains $k$, so has dimension at least one. Hence
\[ p_a(Z)=1-h^0(\mathcal O_Z)+h^1(\mathcal O_Z)\le h^1(\mathcal O_X)=p_a(X). \]
No reducedness of $X$ or $Z$ is needed. \hfill$\square$

\exercise{5.3}{For proper integral $X/k$ and a Cartier divisor of positive degree, prove ampleness using normalization $\pi:X'\to X$. (a) For torsion-free coherent $\mathcal F$, obtain $0\to\pi_*\pi^!\mathcal F\to\mathcal F\to\mathcal H\to0$ with finite-support $\mathcal H$. (b) Prove the required $H^1$ vanishing upstairs and the projection formula. (c) Descend vanishing and conclude.}
\pf
(a) The quotient $\mathcal S$ in $0\to\mathcal O_X\to\pi_*\mathcal O_{X'}\to\mathcal S\to0$ has finite support. A map from $\mathcal S$ to torsion-free $\mathcal F$ is zero. Applying sheaf Hom gives an injection $\mathcal H om_X(\pi_*\mathcal O_{X'},\mathcal F)\to\mathcal F$. On the normal locus it is an isomorphism, so its cokernel has finite support. Finite duality defines the coherent $\mathcal O_{X'}$-module $\pi^!\mathcal F$ by precisely this Hom module.

(b) The normal proper curve $X'$ is projective, and $\deg\pi^*D=\deg D>0$. Riemann--Roch or the positive-degree ampleness criterion on the normal projective curve makes $\pi^*\mathcal O_X(D)$ ample. Thus Serre vanishing gives $H^1(X',\pi^!\mathcal F\otimes\pi^*\mathcal O_X(nD))=0$ for all sufficiently large $n$. The projection formula holds because the twisting sheaf is invertible; on affine charts it is the tensor-product identification for a finite algebra.

(c) Tensor (a) by $\mathcal O_X(nD)$. Finite morphisms have no higher direct images on quasi-coherent sheaves; (b) therefore makes the $H^1$ of the first term zero. The cokernel has finite support and no $H^1$, so $H^1(X,\mathcal F(nD))=0$. For arbitrary coherent $\mathcal F$, first remove its torsion subsheaf, which has finite support and no positive cohomology. This gives the same vanishing. Higher cohomology already vanishes by dimension, and Serre's cohomological criterion for ampleness proves the claim. \hfill$\square$

\exercise{5.4}{For any proper curve $X/k$, (a) construct an effective Cartier divisor meeting every irreducible component; (b) prove it ample and deduce projectivity.}
\pf
There are finitely many components and associated points. On each component choose a closed point outside the other components and the embedded associated points. Such points exist by density of closed points. At the chosen point the one-dimensional local ring has no embedded associated prime, so prime avoidance gives a nonzerodivisor in its maximal ideal. Extend it to a function on an affine neighborhood, shrink to make it regular as an element and to remove all unwanted points of its zero set. Its remaining zero set is the chosen point. Glue its local principal divisor with the zero divisor on the complement. This gives an effective Cartier divisor at that point; summing gives $D$ meeting every component.

On each reduced integral component the restriction has positive degree and is ample by 5.3. Ampleness of an invertible sheaf can be checked on the finitely many reduced components: first pass through the finite normalization and use descent of ampleness for finite surjective maps, then use invariance under nilpotent thickening (equivalently filter coherent sheaves by powers of the nilradical in Serre's criterion). Thus $D$ is ample on $X$. A proper finite type scheme over a field with an ample invertible sheaf is projective. \hfill$\square$

\exercise{5.5}{If $X/k$ is a separated curve and no irreducible component is proper, show $X$ affine.}
\pf
First let $X$ be integral. Its normalization $X'$ is finite over $X$ and is a nonproper normal separated curve. Its smooth/regular projective completion has a nonempty finite boundary, so $X'$ is affine: that boundary is a positive-degree divisor on the normal projective curve, and the complement of an effective ample divisor is affine.

For torsion-free coherent $\mathcal F$ the construction in 5.3(a), valid without properness, injects $\pi_*\pi^!\mathcal F$ into $\mathcal F$ with finite-support cokernel. The first term has zero $H^1$ because $X'$ is affine; the cokernel also does, so $H^1(X,\mathcal F)=0$. Torsion sheaves have finite support, so the same holds for every coherent sheaf. The argument works for reduced reducible $X$ by taking the disjoint union of its component normalizations: every component normalization is affine, and the kernel and cokernel away from the generic points have finite support. For a nonreduced curve use the nilradical filtration and affineness of the reduction. Serre's criterion now gives affineness of $X$; equivalently affineness is invariant under nilpotent thickenings. Separatedness is used for the normal projective completion and for coherent cohomology on these curves. \hfill$\square$

\exercise{5.6}{For a reduced curve over algebraically closed $k$, let $n_x$ be the number of global irreducible components through $x$. Prove $\delta_x\ge n_x-1$, with equality only for an ordinary multiple point.}
\pf
Let $B=(\pi_*\mathcal O_{X'})_x$ and $A=\mathcal O_{X,x}$, and let $m$ be the number of points of the normalization over $x$, that is, the number of branches. Each component through $x$ has at least one branch, so $m\ge n_x$. Reduction of $B$ at its maximal ideals gives a surjection $B\to k^m$, while elements of $A$ have diagonal residue. Thus $B/A$ surjects onto $k^m/k$, giving
\[ \delta_x=\dim_k(B/A)\ge m-1\ge n_x-1. \]
If the endpoints are equal then $m=n_x$ and $\delta_x=m-1$, precisely Liu's definition of an ordinary multiple point. Equivalently $A$ is the subring of $B$ of tuples with equal residues. \hfill$\square$

\exercise{5.7}{For $F(y,z)=\sum F_i(y,z)$ over algebraically closed $k$, $F_0=0$, examine $X=V(F)$ at the origin. (a) Show smoothness equivalent to $F_1\ne0$; when $F_1=0,F_2\ne0$, reduce $F_2$ to $yz$ or $y^2$. (b) In the split case construct formal coordinates with $F=uv$. (c) In the double-tangent case write $F=y^2(1+\epsilon)+z^{r_1}H_1(z)y+z^{r_2}H_2(z)$, with $H_i(0)\ne0$, and put $r=\min(r_1,\lfloor r_2/2\rfloor)$. Prove $w=y/z^r$ integral and assess the asserted $\delta=r$. (d),(e) Assess the proposed branch-count criteria from the residual discriminant and from $r_1>r$ in characteristic two.}
\pf
(a) The hypersurface Jacobian criterion at a rational point says exactly that the linear term is nonzero. A nonzero binary quadratic over an algebraically closed field factors into linear factors; independent factors give $yz$ after a linear change, and a repeated factor gives $y^2$, absorbing a nonzero scalar.

(b) Start with $y_1=y,z_1=z$. Suppose $F-y_nz_n\in(y,z)^{n+2}$. Its leading homogeneous error can be written $yB+zA$ with $A,B$ homogeneous of degree $n+1$. Set $y_{n+1}=y_n+A$, $z_{n+1}=z_n+B$, with signs chosen to cancel that error. The difference between $y_n,z_n$ and $y,z$ only affects the next order, and $AB\in(y,z)^{n+3}$. Thus $F-y_{n+1}z_{n+1}\in(y,z)^{n+3}$. The sequences converge to $u,v$ with identity linear terms, hence define a formal coordinate automorphism. It sends $F$ to $uv$, giving $\widehat{\mathcal O}_{X,0}=k[[u,v]]/(uv)$ and a node.

(c) Assume the curve is reduced near the origin, as needed for the normalization invariants. The variable $z$ is a nonzerodivisor, since it does not divide $F$. Dividing the equation by $z^{2r}(1+\epsilon)$ gives
\[ w^2+\frac{z^{r_1-r}H_1(z)}{1+\epsilon}\,w+
\frac{z^{r_2-2r}H_2(z)}{1+\epsilon}=0. \]
All coefficients lie in the local ring, and the exponents are nonnegative, so $w$ is integral.

\textbf{The formula $\delta=r$ is false in general.} In characteristic zero take
\[ F=(y-z^2)^2-z^6=(y-z^2-z^3)(y-z^2+z^3). \]
Here $r_1=2,r_2=4,r=2$, while there are two smooth branches with intersection multiplicity three. Hence $\delta=0+0+3=3$, not $2$. Explicitly the difference between their graph functions is $2z^3$, and the quotient measuring their failure to glue independently has length three.

There is nevertheless a useful lower bound. By Weierstrass preparation, the completed ring is a free rank-two $R=k[[z]]$-module with basis $1,y$. Its normalization contains the lattice $R+Rw$, by the integrality just proved. This lattice modulo $R+Ry$ has length $r$. Thus $\delta\ge r$, with equality exactly when this lattice is already the normalization. For a monic quadratic equation the transformed equation in $(w,z)$ computes the extra normalization defect; one must not assume that transformed curve nonsingular.

(d) In the characteristic-zero example the printed residual discriminant is $4z^2$, which vanishes at zero, yet the number of branches is two, not one. The valid implication is that a \emph{nonzero} residual discriminant gives two distinct residual roots and hence two smooth branches of the transformed curve, by Hensel's lemma. A zero residual discriminant requires further analysis and can yield either one or two branches.

(e) In characteristic two take
\[ F=(y+z^2)(y+z^2+z^3)=y^2+z^3y+z^4+z^5. \]
Then $r_1=3,r_2=4,r=2$, so $r_1>r$, but there are two branches, contradicting the printed criterion. Their intersection multiplicity is again three. If $r_1=r$, the residual quadratic in $w$ has nonzero linear coefficient, hence two distinct roots over $k$ and two branches. When $r_1>r$ the repeated residual root alone does not determine the branch count. These counterexamples address the unrestricted assertions, rather than silently applying the formulas only to already normalized transforms. \hfill$\square$

\exercise{5.8}{For an Artinian local $k$-algebra $A$ with residue field $k$, prove $A^*\simeq k^*\times U$ with $U$ unipotent.}
\pf
The structure map splits the residue map on scalars. Every unit is uniquely $c(1+a)$ with $c\in k^*$, $a\in\fm_A$, so $U=1+\fm_A$. The maximal ideal is nilpotent, and the filtration $1+\fm_A^i$ has quotients
\[ (1+\fm_A^i)/(1+\fm_A^{i+1})\simeq(\fm_A^i/\fm_A^{i+1},+), \]
since multiplication modulo the next power is addition. These are finite-dimensional additive groups: $A$ has finite length with all simple factors $k$, hence is finite-dimensional over $k$. This filtration exhibits $U$ as unipotent, both in the elementary sense of the text and as an algebraic group. \hfill$\square$

\exercise{5.9}{Let reduced $X/k$, $k=\bar k$, have exactly one singular point $x$, an ordinary multiple point. If $f:X'\to Z$ on its normalization takes all $y_i\in\pi^{-1}(x)$ to the same point, prove unique descent to $X$.}
\pf
Write $A=\mathcal O_{X,x}$ and $B=(\pi_*\mathcal O_{X'})_x$. Ordinariness and the argument of 5.6 give
\[ A=\{b\in B:b(y_1)=\cdots=b(y_m)\}. \]
Indeed $A$ is contained in the ring on the right, and both its quotient in $B$ and that ring's quotient in $B$ have dimension $m-1$.

Let $z=f(y_i)$ and choose an affine neighborhood $W$ of $z$ in $Z$. By finiteness of $\pi$, the image of the closed complement of $f^{-1}(W)$ is closed and misses $x$. Thus after shrinking to a neighborhood $V$ of $x$, $f(\pi^{-1}V)\subseteq W$. Pullbacks of functions on $W$ have equal residues at all $y_i$, so their germs lie in $A$. Choose finitely many algebra generators and shrink $V$ again to make their pullbacks regular on $V$. Their relations hold there because $\mathcal O_V$ injects into its normalization. This defines $V\to W$ lifting $f$. Away from $x$, $\pi$ is an isomorphism, so this map glues with the uniquely forced map on $X\setminus\{x\}$. Uniqueness follows locally from the same injective map of functions, and globally from surjectivity of $\pi$ and agreement on the indicated opens. \hfill$\square$

\exercise{5.10}{Let reduced connected projective $X/\bar k$ be the union of connected closed curves $X_1,\ldots,X_n$ without shared irreducible components. (a) Prove $a(X)=\sum_i a(X_i)$ for Abelian ranks. (b) For $n=2$ prove $t(X)=t(X_1)+t(X_2)+\#(X_1\cap X_2)-1$; deduce $t(X)\ge\sum_i t(X_i)$ in general.}
\pf
The Abelian quotient of the generalized Jacobian is the product of the Jacobians of the normalized irreducible components. The partition into the $X_i$ partitions those same components, so the dimensions add, proving (a). If a subscheme $X_i$ has a nilpotent structure at its intersections, its Abelian and toric ranks are unchanged upon reduction; thus it suffices to use the reduced curves.

Here is a direct count for (b). For a reduced connected curve $Y$, let $c(Y)$ count its irreducible components and let $m_x(Y)$ be its number of branches at $x$. The normalization unit sequence gives
\[ t(Y)=\sum_{x\in Y}(m_x(Y)-1)-c(Y)+1, \]
where only finitely many terms are nonzero. Each local unit quotient contributes a torus of rank $m_x-1$, and constants on the $c(Y)$ normalization components contribute a torus of rank $c(Y)-1$ to be divided out; unipotent local factors do not alter this rank.

At a point meeting $q_x$ of the subcurves $X_i$, branches of $X$ are the disjoint union of their branches, so
\[ (m_x(X)-1)-\sum_{i:x\in X_i}(m_x(X_i)-1)=q_x-1. \]
Since $c(X)=\sum_i c(X_i)$ and all the curves in question are connected,
\[ t(X)-\sum_i t(X_i)=\sum_x(q_x-1)-n+1. \]
For two subcurves this is exactly $\#(X_1\cap X_2)-1$. In general form the finite incidence graph with one vertex for each subcurve and one for each intersection point, joining a point to the subcurves containing it. Connectedness of $X$ makes this graph connected. Its first Betti number is $\sum_x(q_x-1)-n+1$, a nonnegative integer, proving the inequality. \hfill$\square$

\end{document}
